\answer{ex:09:1}{(a)~$P_\infty=0.2$, mémoire $20\,\text{s}$. (b)~$P_\infty=1$, mémoire $1\,\text{s}$. (c)~La mémoire est proportionnelle à l'amplitude du bruit: six fois.}
\answer{ex:09:2}{$P(t)=P_\infty\coth(t\sqrt{q/R})$. (a)~$\frac12\sqrt{R/q}=5\,\text{s}$, deux fois moins que la mémoire. (b)~$t=\frac12\ln21\,\sqrt{R/q}\approx15\,\text{s}$: c'est le temps pendant lequel \nt{ACET} doit rester élevé.}
\answer{ex:09:3}{Variance $qT/2+R/(2T)$; $T^*=\sqrt{R/q}$, variance $\sqrt{qR}=P_\infty$, comme pour le filtre~\eqref{eq:kalman}; croissance d'un facteur $(\kappa+1/\kappa)/2$: $1.25$ et $5.05$.}
\answer{ex:09:4}{$a>a_w=1-\theta/(mw)$: $0.15$ pour $w=0.041$, $0.22$ pour $w=0.045$.}
\answer{ex:09:5}{Les erreurs apparaissent vers $M\approx40$--$60$; pour $M=80$, $1.9$ neurone superflu; pour $M=100$, part du motif $0.86$; l'interférence des autres motifs a une variance $\approx(M-1)p(1-p)/N$, d'où $M_{\max}\approx80$, $M_{\max}\propto N/p$.}
\answer{ex:09:6}{Par exemple $\eta(a)=\eta\min\bigl(1,\max(0,(a-0.3)/0.6)\bigr)$. Avec $\eta a$, les trois neurones étrangers sont entrés dans le motif; avec $\eta(a)$, aucun, et l'écriture pour $a\ge0.9$ est la même (part $1.00$).}
\answer{ex:09:7}{(a)~$40$ répétitions; avec la fonction $\eta(a)$ de l'exercice~\ref{ex:09:6}, jamais. (b)~$200$ répétitions, $10$ nuits.}
